\ifx\danslelivre\undefined
\documentclass[10pt]{fichiers-configuration-latex/smfart}

\begin{document}

\begin{center}
Commentaires sur SGA~4-I\\
P. Dehornoy
\end{center}


\else
\chapter{Commentaires sur SGA~4-I}
%\addtocontents{toc}{par P. Dehornoy}
\begin{center}
par P. Dehornoy
\end{center}
\fi

\begin{enumerate}
\item Tout comme le livre I de Bourbaki, l'exposé SGA4-1 est le
reflet d'une vue extrêmement naïve de la théorie des
ensembles telle qu'on pouvait à la rigueur la garder avant
les théorèmes d'incomplétude de Gödel. Les ensembles sont
vus essentiellement comme des objets algébriques
effectivement définissables, et la théorie des ensembles est
conçue comme la manipulation d'expressions formelles
compliquées à base d'ensemble vide et d'accolades
innombrables, sans aucun ingrédient conceptuel
intrinsèquement logique et infinitiste. La description de la
structure $U_1$ la structure formée par les ensembles
usuellement appelés héréditairement finis et la remarque que
cet objet est presque un univers à la mineure exception
qu'il ne contient pas d'ensemble infini sont typiques:
l'introduction d'ensembles infinis est cruciale, et c'est
exactement le point où la théorie des ensembles commence,
par opposition à l'arithmétique, qui est le monde où $U_1$
vit naturellement. Le point le plus négatif dans cette
approche algébrique de la théorie des ensembles est la
confusion du monde des formules logiques et celui des
ensembles; dans toutes les approches modernes, ces mondes
sont distincts, les formules sont des objets syntaxiques
(mots), alors que les ensembles sont des objets sémantiques
(objets mathématiques quelconques) sur lesquelles les
formules opèrent. Cette distinction --- dont on comprend
bien que Bourbaki, dans sa démarche unificatrice, ait
cherché à l'éviter, mais c'était une option malheureuse ---
est la seule solution pour parvenir à la notion de modèle
abstrait de la théorie des ensembles (mêmes formules, mais
différents ensembles), ingrédient essentiel de tous les
développements ultérieurs de la théorie des ensembles, en
particulier les théorèmes d'indépendance relative de Cohen
et les récents résultats laissant espérer une solution
prochaine du problème du continu. On ne saurait donc trop
insister sur le fait que l'axiome « il existe un ensemble
infini » n'est pas juste un petit axiome supplémentaire, mais
est l'axiome essentiel, et, en particulier, qu'il ne saurait
être déduit des autres axiomes comme le démontre par exemple
le second théorème d'incomplétude de Gödel puisqu'en
montrant l'existence de l'ensemble $U_1$ la théorie des
ensembles prouve la non-contradiction de l'arithmétique.
L'un des aspects les plus regrettables dans l'approche
algébrique de la théorie des ensembles est d'induire le
sentiment erroné que la théorie des ensembles n'est qu'une
sorte de prolongement infini plus ou moins trivial de la
théorie des ensembles finis, masquant ainsi les problèmes
véritables et suscitant de ce fait un compréhensible
désintérêt pour tous les développements ultérieurs.

\item Cela dit, le point de vue de SGA diffère de celui du
traité de Bourbaki en cela que la théorie des ensembles n'y
est pas invoquée comme base d'une théorie du grand tout,
mais simplement comme contexte cohérent où développer une
théorie spécifique des catégories et des faisceaux. Cette
approche est beaucoup plus modeste, et, de ce fait, les
aspects négatifs signalés plus haut deviennent secondaires.
Au total prévaut une approche pragmatique, dans laquelle il
s'agit simplement de recenser les outils ensemblistes dont
on aura besoin dans la suite.
Si ces choses avaient à être rédigées aujourd'hui, on
pourrait tout à fait reprendre la même approche, et il
serait simplement naturel de la rattacher au système
standard de Zermelo-Fraenkel plutôt qu'à la théorie des
ensembles de Bourbaki, qui, quoique à peu près équivalente,
n'est plus guère utilisée. Il serait peut-être opportun
d'éviter le terme « univers », souvent utilisé pour désigner
un modèle de la théorie ZF, c'est-à-dire une structure
$(M,E)$ quelconque dans laquelle les axiomes de ZF sont
satisfaits -- mais ce point de vocabulaire est mineur.

\item Comme chacun sait, il serait contradictoire de considérer
un ensemble de tous les ensembles (paradoxe de Russell), ce
qui entraîne immédiatement des problèmes pour introduire des
catégories dans un contexte de théorie des ensembles. Il
faut insister sur le fait que ces problèmes de taille des
ensembles  (par opposition aux problèmes d'existence et de
complexité) sont, d'un point de vue de théorie des
ensembles, des problèmes mineurs, pour des raisons précises
qui vont être expliquées ci-dessous -- et, de là, qu'il n'y
a aucune raison sérieuse d'opposer les points de vue
ensembles et catégoriques. Le point de vue adopté dans
SGA4-1 paraît excellent, et, de fait, il semble le seul pour
pouvoir développer sans problème toutes les constructions
basées sur des itérations. Une solution standard pour éviter
les problèmes d'ensemble de tous les ensembles consiste à
introduire des classes comme second type d'objet, de sorte
qu'il existe une classe de tous les ensembles, mais qu'on ne
forme pas de classe de classes en général.
Cette solution, qui est pertinente si on veut considérer une
catégorie particulière, serait peu adaptée dès lors qu'on
veut pouvoir construire en toute généralité des catégories
de catégories, catégories dérivées, etc. Il est donc
précieux de pouvoir considérer les objets et les morphismes
d'une catégorie quelconque comme des ensembles. Comme cela
ne peut être de vrais ensembles à cause du paradoxe de
Russell, on introduit une notion relative d'ensemble, les
$U$-ensembles, ou ensembles selon $U$, où $U$ est une
approximation convenable des ensembles, à savoir une qui est
close par suffisamment d'opérations ensemblistes.

\item D'un point de vue moderne, on peut remarquer qu'il y a
certainement trop de généralité dans les univers du SGA4-1.
De fait, dans le système standard ZF, les ensembles sont
munis d'une filtration par les ordinaux appelée rang. La
notation traditionnelle est $V_\alpha$ pour les ensembles de
rang moindre que $\alpha$ (« V » comme « von Neumann »). Par
exemple, l'univers $U_0$ est $V_0$, et $U_1$ est $V_\omega$.
Pour tout ordinal limite $\lambda$ (c'est-à-dire un ordinal
qui n'est pas de la forme $\alpha + 1$), l'ensemble
$V_\lambda$ est presque un univers, et, d'un autre côté,
c'est un modèle de la théorie des ensembles de Zermelo (qui
est Zermelo-Fraenkel moins la clôture par fonction). La
remarque est qu'il me semble que les seuls univers utiles
pour SGA4-1 sont des $V_\lambda$. Dans l'approche
bourbachique, cela peut paraître une faible restriction; en
fait, dans la jungle de la théorie des ensembles, il existe
une immense variété d'univers au sens de SGA4-1, et les
$V_\lambda$ en sont des exemples fort particuliers. Ceci
traduit le fait que le seul problème posé concerne la taille
des ensembles, et qu'une filtration par le rang est
suffisante pour l'éviter. 

\item Le contexte naturel aujourd'hui serait donc plutôt la
notion de $V_\lambda$-catégorie. Pour assurer que l'ensemble
$V_\lambda$ soit clos par suite et fonction, il faut
réclamer davantage que le fait que l'ordinal $\lambda$ soit
un ordinal limite. Une condition suffisante, et très
probablement nécessaire, est que $V_\lambda$ soit un modèle
de ZF. On est donc conduit à supposer l'existence d'une
classe non bornée d'ordinaux $\lambda$ tels que $V_\lambda$
soit modèle de ZF. En vertu du second théorème
d'incomplétude de Gödel, l'existence de tels ordinaux, et
même d'un seul tel ordinal, est une hypothèse indémontrable
à partir du système standard ZF, toujours pour la même
raison que cette existence démontre la non-contradiction de
ZF et donc ne peut être démontrée à partir de ZF -- ceci n'a
rien à voir avec Cohen, ainsi que le texte de SGA le suggère
prudemment.
Par contre, comme SGA le remarque justement, une condition
suffisante pour que $V_\lambda$ soit un modèle de ZF -- et
donc un univers au sens de SGA -- est que $\lambda$ soit un
cardinal (fortement) inaccessible, et, par conséquent, un
système axiomatique suffisant pour justifier toutes les
constructions ultérieures est le système « ZF + il existe une
classe non bornée de cardinaux inaccessibles ».

\item La situation est donc la suivante, que le système
axiomatique à partir duquel les constructions sont
effectuées est le système « ZF + il existe une classe non
bornée de cardinaux inaccessibles », qu'on notera ZFI, et non
le système ZF standard. Les questions naturelles à ce stade
sont: (i) Le système ZFI est-il raisonnable? (ii) Le système
ZFI est-il nécessaire? 

\item La question (i) se divise elle-même en deux questions
distinctes: (i') Le système ZFI est-il exempt de
contradictions? (i'') Le système ZFI est-il objet d'un
consensus? Pour (i'), la réponse est claire: en vertu du
second théorème d'incomplétude de Gödel, il ne peut exister
de preuve de la cohérence de ZFI, ni de preuve de la
cohérence relative de ZFI par rapport à ZF. Le système ZFI
est un système d'axiomes strictement plus fort que ZF, et,
de même qu'il se pourrait que l'arithmétique soit cohérente
alors que ZF est contradictoire, il se pourrait que ZF soit
cohérent et que ZFI soit contradictoire. Cela dit, aucun
résultat à ce jour ne suggère une telle situation, et, à
côté de la question mathématique (i'), se pose la question
sociologique (i''), et c'est à ce niveau qu'il peut être
intéressant d'observer les évolutions récentes. A l'époque
où SGA a été rédigé, l'existence de cardinaux inaccessibles,
et plus généralement de ce qu'on commençait à appeler des
grands cardinaux, était considéré comme une hypothèse
technique parmi d'autres, sans qu'il existe aucun consensus
pour considérer cette hypothèse comme vraie, c'est-à-dire
pour proposer de l'inclure systématiquement dans les axiomes
de base. Quarante ans plus tard, la vision des experts a
considérablement évolué, et les axiomes de grands cardinaux
sont devenus une sorte de toile de fond largement acceptée.
Le sentiment général est, d'une part, que le système ZF
n'épuise absolument pas notre intuition de la notion
d'ensemble et que, par conséquent, il s'agit de poursuivre
l'étude de façon à dégager des propriétés supplémentaires
qu'un consensus propose de prendre comme point de départ,
et, d'autre part, que les axiomes de grands cardinaux, donc
en particulier l'existence d'une classe non bornée de
cardinaux inaccessibles qui est l'un des plus modérés
d'entre eux, font partie de ces propriétés supplémentaires
qu'il est légitime d'inclure dans les axiomes. Cette opinion
ne procède pas d'une évidence a priori ou de discutables
expériences de pensée, mais d'un faisceau convergent de
théorèmes montrant que la description du monde mathématique
résultant de l'introduction de ces axiomes est conforme à ce
que l'intuition recommande. Typiquement, on sait depuis les
résultats de Steel--Martin et Woodin dans les années 1980
que le système obtenu en ajoutant à ZF l'axiome dit de
détermination projective -- qui est un axiome de grand
cardinal (beaucoup) plus fort que l'existence d'une classe
non bornée de cardinaux inaccessibles -- est, en un sens
mathématiquement précis, la bonne axiomatisation de l'infini
dénombrable à la façon dont ZF est la bonne axiomatisation
du fini. Dans ce contexte, le fait que SGA parte du système
ZFI au lieu de ZF apparaît donc aujourd'hui comme
parfaitement inoffensif et légitime, bien davantage qu'il y
a quarante ans.

\item Venons-en à la question (ii), c'est-à-dire à la nécessité
d'utiliser ZFI plutôt que ZF comme base axiomatique. Là
encore, le problème est double, suivant qu'il s'agit
d'établir des résultats généraux du type « pour toute
catégorie, il existe... » ou des résultats particuliers du
type « pour telle catégorie particulière, il existe... » ou
reposant sur des énoncés de ce type. Dans tous les cas, il
n'est pas certain qu'on ait besoin du fait que tout ensemble
appartient à un univers, pour lequel ZFI est condition
suffisante: il se peut par exemple qu'on n'utilise que des
suites dénombrable d'itérations, et qu'alors une condition
suffisante soit l'existence d'une suite dénombrable de
cardinaux inaccessibles, mais ZFI semble bien la plus simple
des hypothèses susceptibles de marcher dans tous les cas
(l'existence d'une tour de $V_\lambda$ modèles de ZF est, du
point de vue logique, une hypothèse plus faible, mais moins
intuitive). Par contre, si on considère une catégorie
particulière, il est plus que probable qu'on n'a pas besoin
d'une hypothèse aussi forte que ZFI. En particulier, il
existe un phénomène général dit de réflexion montrant que,
pour tout énoncé ensembliste $E$, il existe un ordinal
$\lambda$ tel que $E$ est équivalent à l'énoncé obtenu en
bornant toutes les quantifications à $V_\lambda$. Ceci ne
prouve rien en soi, mais ce genre de phénomène suggère qu'on
n'utilise probablement jamais la pleine hypothèse qu'un
ensemble appartient à un univers (au sens de SGA), mais
simplement qu'il appartient à un ensemble $V_\lambda$ qui
satisfait non pas tous les axiomes de ZF, mais simplement
suffisamment de ces axiomes pour les besoins spécifiques,
typiquement un nombre fini de tels axiomes (il est important
de réaliser que ZF est un système d'axiomes infini, tout
comme le système de Bourbaki à cause de la quantification
sur les termes en $\tau$). Or, pour chaque fragment fini de
ZF, l'existence d'un tel $V_\lambda$ se démontre dans ZF,
sans qu'il soit besoin de faire appel à un cardinal
inaccessible. Il apparaît donc probable -- sans que je
puisse l'affirmer, il faudrait vérifier pas à pas, et les
remarques du point 7 montrent que le jeu n'en vaut peut-être
pas la chandelle -- que, lorsqu'il s'agit de démontrer un
théorème donné, par exemple, Fermat, en utilisant pour cela
un certain nombre de résultats sur des catégories, la
démonstration puisse se miniaturiser de façon à n'utiliser à
chaque fois que des univers partiels suffisants et, de ce
fait, puisse être entièrement formalisée à partir de ZF --
voire même d'un système significativement plus faible du
point de vue logique ainsi que c'est le cas de la plupart
des résultats mathématiques usuels. Un parallèle peut-être
éclairant peut être proposé avec l'énoncé « Tout corps
possède une clôture algébrique ». Si on souhaite l'énoncé
général, le seul cadre connu utilise le lemme de Zorn, donc
l'axiome du choix. Par contre, si on souhaite seulement
utiliser le fait que les rationnels, ou tel ou tel corps
explicitement énuméré, admettent une clôture algébrique,
alors il suffit d'utiliser une construction inductive qui
n'utilise aucun choix. La seule différence est qu'on sait
(depuis 1938) que l'axiome du choix, même non impliqué par
les axiomes de ZF, est inoffensif au sens où son adjonction
ne saurait introduire de contradiction, alors que
l'existence d'une classe non bornée de cardinaux
inaccessibles n'est pas inoffensive -- mais on a dit en 7
qu'elle est néanmoins l'objet d'un consensus large.

\item Dernière remarque: Les développements de l'appendice sur
les ensembles appelés artiniens (usuellement appelés
ensembles purs, ou bien fondés) et la possibilité de se
restreindre à de tels ensembles correspondent exactement à
l'inclusion de l'axiome dit de fondation dans le système de
Zermelo-Fraenkel, lequel correspond aussi à l'affirmation
ci-dessus utilisée que tout ensemble a un rang. Cette
hypothèse est donc implicite en théorie des ensembles
standard.
\end{enumerate}

\ifx\danslelivre\undefined
\end{document}
\fi
